\section{Labeled graphs and the parity criterion}\label{sec:algebra}

The algebraic part of the construction is a parity calculation on finite
graphs.  We first isolate this calculation from the geometric work needed
to ensure that the resulting group is torsion-free and that a specified
coefficient does not vanish.  The use of immersed graphs and their cones
follows the construction in \cite[Proposition~5.1 and Section~6]{OpenAI2026ZD}.

A \emph{signed alphabet} is a finite set $T$ with a fixed-point-free
involution $t\mapsto\bar t$.  Let $F$ be the rose with one vertex $*$ and
one geometric edge for each pair $\{t,\bar t\}$; its two orientations
carry the labels $t$ and $\bar t$.  A $T$-labeled graph $\Gamma$ has a
label on every oriented edge, with opposite orientations carrying inverse
labels.  It therefore has a label map $\Gamma\to F$ taking each edge
homeomorphically onto the appropriate edge of the rose.  The map is an
\emph{immersion} if the labels on the edges leaving any one vertex are
distinct.  Write $S_x\subseteq T$ for these outgoing labels.  For $t\in
S_x$, denote by $x\cdot t$ the endpoint of the unique edge leaving $x$
with label $t$.

The length of an edge path is its number of edge traversals, counted with
multiplicity.  A path is \emph{immersed} if it has no consecutive edge
and reverse edge.  A closed path is \emph{cyclically immersed} if this
also holds at the closing turn.  For a graph immersed in $F$, these
conditions say precisely that the label word is freely reduced, or
cyclically reduced, respectively.  We multiply labels in the order in
which a path traverses its edges.

Let $\Gamma_A$ and $\Gamma_B$ be finite graphs immersed in $F$, with
vertex sets $A$ and $B$, and choose roots $x_A\in A$ and $x_B\in B$.
Put $\Gamma=\Gamma_A\sqcup\Gamma_B$.  For each connected component
$\Lambda$ of $\Gamma$, let $K_\Lambda$ be its abstract cone, and form
\begin{equation}\label{eq:cone-complex}
 X=F\mathop{\cup}_{\Lambda\to F}
       \bigcup_{\Lambda\in\pi_0(\Gamma)}K_\Lambda,
 \qquad G=\pi_1(X,*).
\end{equation}
Here the base of each cone is attached by its label map.  In particular,
the cone need not embed in $X$.  This is a finite CW complex of dimension
at most two.  Every closed path in $\Gamma$ has trivial label in $G$,
because it bounds in its cone.

Let $A'\subseteq A$ and $B'\subseteq B$ be the vertex sets of the
components containing the respective roots.  For $x\in A'$, let $g_x$
be the label in $G$ of any path from $x_A$ to $x$; define $h_y$ for
$y\in B'$ in the same way.  These elements are independent of the
chosen paths: the concatenation of one such path with the reverse of
another is a closed path in the same component.  Thus
\begin{equation}\label{eq:path-step}
 g_{x_A}=h_{x_B}=1,\qquad
 g_{x\cdot t}=g_x t,\qquad h_{y\cdot t}=h_y t,
\end{equation}
whenever the indicated step exists, where a letter also denotes its
image in $G$.

\begin{proposition}[Parity criterion]\label{prop:parity}
Let $\Gamma_A,\Gamma_B$ be finite loopless graphs immersed in the same
rose, with roots and cone group as above.  Suppose that
\begin{equation}\label{eq:parity-assumptions}
 \begin{aligned}
 |S_x\cap S_{x'}|&\equiv1\pmod2
   &&\bigl(x,x'\in A',\ (x,x')\ne(x_A,x_A)\bigr),\\
 |S_{x_A}|&\equiv0\pmod2,\\
 |S_x\cap S_y|&\equiv1\pmod2
   &&\bigl(x\in A',\ y\in B'\bigr).
 \end{aligned}
\end{equation}
Suppose also that the roots are protected in the following sense:
\begin{equation}\label{eq:root-protection}
 \begin{gathered}
 \text{if a path in $\Gamma$ starts at $x_A$ or $x_B$ and ends}\\[-2pt]
 \text{at a different vertex, its label is nontrivial in $G$.}
 \end{gathered}
\end{equation}
Then the finite scalar sums
\begin{equation}\label{eq:scalar-elements}
 a=\sum_{x\in A'}g_x,\qquad
 b=\sum_{x\in A'}g_x^{-1},\qquad
 c=\sum_{y\in B'}h_y^{-1}
 \quad\text{in }\Ftwo[G]
\end{equation}
satisfy $ab=1$, $ac=0$, and $c\ne0$.  Consequently $ba\ne1$.
\end{proposition}

\begin{proof}
On $A'\times A'$ form the simultaneous-step graph: for every common
outgoing label $t$ at $x$ and $x'$, join $(x,x')$ to
$(x\cdot t,x'\cdot t)$.  The step with label $\bar t$ at the latter
vertex is the reverse of this edge, rather than a second edge.
Different steps with the same pair of endpoints remain distinct edges.
The graph is loopless because the factor graphs are loopless, and its
degree at $(x,x')$ is exactly $|S_x\cap S_{x'}|$.

The group element contributed to $ab$ by a vertex is constant on each
component of this graph.  Indeed, by \eqref{eq:path-step},
\[
 g_{x\cdot t}g_{x'\cdot t}^{-1}
   =(g_xt)(g_{x'}t)^{-1}=g_xg_{x'}^{-1}.
\]
In every finite graph the number of vertices of odd degree is even,
because the sum of the degrees is twice the number of edges.  By
\eqref{eq:parity-assumptions}, every component that does not contain
$(x_A,x_A)$ therefore has an even number of vertices.  Its contributions
to $ab$ cancel in characteristic two.  In the component containing
$(x_A,x_A)$, that vertex is the unique vertex of even degree, so the
component has odd cardinality.  Its constant contribution is
$g_{x_A}g_{x_A}^{-1}=1$.  This proves $ab=1$.

Apply the same construction to $A'\times B'$.  Its degree at $(x,y)$ is
$|S_x\cap S_y|$, which is always odd.  Hence every component has even
cardinality.  The contribution $g_xh_y^{-1}$ to $ac$ is constant along
each simultaneous step by the same calculation, so $ac=0$.

Finally, $h_{x_B}=1$, and \eqref{eq:root-protection} gives $h_y\ne1$
for every $y\in B'\setminus\{x_B\}$.  Thus the coefficient of the
identity in $c$ is exactly one, and $c\ne0$.  If $ba=1$, associativity
would give $c=(ba)c=b(ac)=0$, a contradiction.
\end{proof}

Only protection of $x_B$ is needed for the last coefficient calculation.
We retain the symmetric condition \eqref{eq:root-protection} because
the geometric argument will establish it for both roots.  The remaining
task is to construct finite immersed graphs satisfying
\eqref{eq:parity-assumptions} whose cone complex is aspherical and whose
roots satisfy \eqref{eq:root-protection}.
